\chapter{Introduction}

We study whether interpretable macroscopic PDE models can be identified from noisy interacting-agent trajectories, and whether the discovered equations reflect the distinct physical mechanisms of each dynamical regime. The thesis pairs WSINDy~\cite{brunton2016sindy,messenger2021wsindy} with two POD-based~\cite{berkooz1993pod} latent forecasters that have deliberately different inductive biases: a linear multivariate autoregression (MVAR) and a nonlinear long short-term memory (LSTM) network. Starting from trajectories $\{x_i(t)\}_{i=1}^N \subset \Omega \subset \mathbb{R}^2$ generated by alignment and attraction--repulsion interactions, we coarse-grain to a smooth density field $\rho(x,t)$ and seek an effective coarse-grained evolution law that is both predictive and physically interpretable. At the macroscopic level, density evolution obeys the conservation law
\begin{equation}
  \partial_t \rho + \nabla\cdot \mathbf{j} = 0,
  \label{eq:conservation_law}
\end{equation}
where $\mathbf{j}(x,t)$ is the agent flux. The modeling task is to identify an effective parameterised closure $\mathbf{j}_{\boldsymbol{\theta}}[\rho,\mathbf{p},\Phi]$ from candidate mechanisms so that
\begin{equation}
  \partial_t \rho(x,t) + \nabla\cdot \mathbf{j}_{\boldsymbol{\theta}}[\rho,\mathbf{p},\Phi](x,t) = 0
  \label{eq:parametrized_flux}
\end{equation}
holds on the observed data, where $\mathbf{p}$ is the polarisation flux and $\Phi=W\ast\rho$ the Morse interaction potential (both defined in Chapter~\ref{ch:background}).  The multi-field dependence reflects the fact that the scalar density equation is not expected to close on its own. In practice, we work in weak form to avoid unstable pointwise differentiation of noisy density fields.

We develop two complementary pipelines. Both start from the same upstream data: particle simulations, KDE-based macroscopic fields, and regime-specific train/test splits. The first is an equation-free reduced-order model (EF-ROM)~\cite{alvarez2025eqfree} that compresses shift-aligned density snapshots via POD, evolves the latent coordinates with a learned time-stepper (MVAR or LSTM), and lifts back to the physical field.  Because coherent flock translation inflates the effective rank under standard POD~\cite{peherstorfer2022kolmogorov,reiss2018shifted_pod,black2020projection,krah2024robust_spod,budanur2015slices,siminos2015periodic,sonday2011alignment}, we register each frame to a common reference via FFT-based cross-correlation on the periodic torus, yielding integer-pixel shifts that remove translational drift without interpolation.  Comparing MVAR and LSTM reveals whether the latent dynamics induced by Vicsek~\cite{vicsek1995vicsek} alignment and D'Orsogna~\cite{dorsogna2006morse} attraction--repulsion remain effectively linear after reduction.

The second pipeline targets interpretability via weak-form sparse identification (WSINDy). Here the learning object differs from the EF-ROM branch: WSINDy does not operate on the POD latent dataset, but instead remains in physical field space and identifies sparse continuous-time equations directly from weak-form systems built on the raw density field and auxiliary multi-fields. We formulate a coupled multi-field system: density $\rho$, polarization fluxes $p_x, p_y$ (from velocity-weighted KDE), and a Morse interaction potential $\Phi = W\ast\rho$ (via FFT convolution). The WSINDy library includes physically motivated candidate terms (continuity, diffusion, Morse forcing, nonlinear diffusion, alignment self-saturation, and density-driven pressure), yielding a three-equation discovery problem for $(\rho_t,\, p_{x,t},\, p_{y,t})$. Compactly supported polynomial test functions and integration by parts replace pointwise derivatives; sparsity is enforced via modified sequential thresholded least squares (MSTLS)~\cite{minor2025mstls} with a dominant-balance criterion~\cite{callaham2021dominant} and automatic test-function width selection. Discovered PDEs are integrated forward using ETDRK4~\cite{cox2002etdrk,kassam2005etdrk4}, which treats stiff linear terms exactly in Fourier space, as a targeted closure stress test on selected identified equations rather than as the primary evaluation.

The central question is not whether the discovered PDEs can ``beat'' the EF-ROM as a forecaster; the two branches answer different questions.  The EF-ROM branch (Part~I, Chapters~\ref{ch:background}--\ref{ch:rom_results}) asks \emph{what is forecastable} from the coarse-grained density movies; the WSINDy branch (Part~II, Chapters~\ref{ch:wsindy_methodology}--\ref{ch:wsindy_results}) asks \emph{what effective mechanisms are identifiable} from the same data.  These are related but not identical questions: a model can forecast accurately without revealing mechanism, and a correctly identified PDE can fail at autonomous rollout because its library does not close the effective law.  Throughout, WSINDy is used \emph{primarily as an identifier:} the closure problem for mean-field particle PDEs (Section~\ref{sec:closure_problem}) means that even a correctly identified sparse PDE may be unclosed, so direct forecast comparisons would be confounded by closure error rather than identification quality.  MVAR and LSTM serve as the forecasting benchmarks; WSINDy contributes complementary mechanistic information about which library terms are active in each dynamical regime.

The central empirical finding is a \emph{closure gap}.  Throughout, we use `closure gap' to mean that the identified PDE, while matching the data in weak form, lacks the mathematical ingredients needed to close the mean-field hierarchy and remain bounded under autonomous integration.  WSINDy recovers sparse, regime-consistent PDE structures with density weak-form $R^2$ exceeding~$0.99$ in eight of nine regimes, yet the identified equations frequently cannot be stably integrated forward in time.  In attractive regimes, the discovered Morse aggregation term creates a positive feedback loop (attraction increases local density, which strengthens attraction) that the candidate library cannot arrest because it lacks a saturation mechanism corresponding to the particle system's short-range repulsion.  This decoupling of identification quality from rollout quality is itself a diagnostic: it reveals which closure ingredients are missing from the effective description at the chosen coarse-graining scale.

Because rollout alone is an insufficient success criterion, we define a composite notion of identification quality.  We use \emph{regime-consistent} throughout to mean that the identified PDE recovers the expected mechanism pattern in a way that is stable, sign-plausible, and well fit in weak form; operationally, the per-regime verdicts in Table~\ref{tab:mechanism_confidence} are based on four observable proxies: continuity-backbone recovery, physically plausible sign structure, cross-scale stability across six support-scale triplets, and strong weak-form fit, with momentum quality used to distinguish strong from moderate or partial outcomes.

\paragraph{Contributions.}
The thesis makes four contributions.

\emph{1.\ Multi-field WSINDy formulation.}
We formulate a coupled weak-form identification problem for the density and polarisation flux equations on finite-$N$ Vicsek--Morse data, incorporating nonlocal Morse interaction features, automatic test-function width selection, and a dominant-balance sparsity criterion.  This extends the scalar WSINDy framework of Messenger and Bortz~\cite{messenger2022meanfield} to a coupled three-field setting with regime-specific libraries.

\emph{2.\ Closure interpretation.}
We show that weak-form identification and autonomous PDE rollout can decouple: WSINDy recovers regime-consistent sparse support patterns and sign structures even where the identified PDE cannot be stably integrated, and we interpret this gap as evidence of closure insufficiency at the selected coarse-graining scale.

\emph{3.\ Shift-aligned POD pipeline.}
We assemble an FFT cross-correlation registration step with $\sqrt{\cdot}$-transformed POD that removes coherent translational drift, enabling faithful low-rank compression of transport-dominated density movies and defining the latent space in which the forecasting benchmark operates.

\emph{4.\ Forecasting benchmark across nine regimes.}
We compare linear (MVAR) and nonlinear (LSTM) latent time-steppers across nine dynamical regimes and four initial-condition families, establishing quantitative forecast ceilings and showing that variable-speed dynamics create a shift-predictability bottleneck that limits both model classes.

\noindent
The thesis does not claim recovery of a unique continuum-limit
PDE independent of coarse-graining; the recovered coefficients are effective
parameters at the chosen data scale, and the closure gap itself is a
central diagnostic.

\paragraph{Software.}The complete implementation accompanying this thesis (the agent simulator, KDE coarse-graining, shift-aligned POD pipeline, MVAR and LSTM latent surrogates, WSINDy discovery engine, ETDRK4 integrator, and all YAML experiment configurations) is released as an open-source Python package at \url{https://github.com/eugenio-macias/wsindy-manifold}. Every figure and table in the following chapters can be reproduced from the scripts and configuration files therein.
A mapping from the mathematical symbols used in the main text to their
corresponding configuration keys and source modules is provided in
Appendix~\ref{app:config_catalogue} (Table~\ref{tab:symbol_config}).

\paragraph{Relation to prior work.}The EF-ROM component follows the restriction--lifting philosophy of Alvarez et al.~\cite{alvarez2025eqfree}; the alignment step draws on shifted POD and symmetry-reduction methods for transport-dominated problems~\cite{sonday2011alignment,budanur2015slices,reiss2018shifted_pod,black2020projection,siminos2015periodic,krah2024robust_spod}, of which we adopt the simplest variant (single-reference FFT cross-correlation on a periodic domain).  The PDE discovery component builds on Rudy et al.'s pointwise sparse regression approach (PDE-FIND)~\cite{rudy2017pdefind} and the subsequent weak-form extension of Messenger and Bortz~\cite{messenger2021wsindy}, which replaces pointwise derivatives with integral test-function projections.  Messenger and Bortz~\cite{messenger2022meanfield} subsequently applied WSINDy to learn scalar mean-field equations from particle data; our contribution extends this to a coupled multi-field identification problem that jointly discovers density and polarisation-flux equations with regime-specific Morse libraries, and systematically applies it across nine regimes.  Concurrent work by Messenger and Bortz~\cite{messenger2025atmospheric} demonstrates that WSINDy can also recover interpretable PDEs for atmospheric phenomena, underscoring the generality of the weak-form identification framework.  The present thesis is also informed by growing interest in geometry-aware modelling, in which macroscopic laws must be consistent with the underlying spatial geometry on which collective motion occurs.  The kinetic-theory literature on Vicsek-type models~\cite{bertin2006boltzmann,degond2008continuum,ihle2011kinetic,toner2005hydrodynamics} provides the external validation baseline used in Section~\ref{subsec:pv_kinetic_validation}.

\paragraph{Reader roadmap.}
Chapter~\ref{ch:background} sets the microscopic modeling and coarse-graining background.  Chapters~\ref{ch:exp_protocol_metrics}--\ref{ch:rom_results} develop and evaluate the EF-ROM forecasting branch.  Chapters~\ref{ch:wsindy_methodology}--\ref{ch:wsindy_results} turn to weak-form PDE identification and closure analysis.  Chapter~\ref{ch:conclusion} synthesises both branches.  Figure~\ref{fig:pipeline_overview} previews the pipeline structure; the headline empirical result, contrasting forecast $R^2$ with weak-form identification $R^2$ across all nine regimes, is presented in Figure~\ref{fig:headline_summary} (Chapter~\ref{ch:wsindy_results}).

\begin{figure}[htbp]
\centering
\definecolor{sharedNavy}{HTML}{2C3E50}
\definecolor{efromTeal}{HTML}{D5F5E3}
\definecolor{efromLabel}{HTML}{117A65}
\definecolor{wsOrange}{HTML}{FDEBD0}
\definecolor{wsLabel}{HTML}{CA6F1E}
\begin{tikzpicture}[
    >=Stealth,
    node distance=0.65cm,
    every node/.style={font=\small},
    shared/.style={
      draw=none, rounded corners=3pt, minimum width=10cm, minimum height=0.7cm,
      fill=sharedNavy, text=white, font=\small\bfseries,
      align=center
    },
    efrom/.style={
      draw=none, rounded corners=3pt, minimum width=5.2cm, minimum height=0.7cm,
      fill=efromTeal, text=black, font=\small, align=center
    },
    wsindy/.style={
      draw=none, rounded corners=3pt, minimum width=5.2cm, minimum height=0.7cm,
      fill=wsOrange, text=black, font=\small, align=center
    },
    branchlabel/.style={font=\normalsize\bfseries\sffamily},
    desc/.style={font=\scriptsize\itshape, text=black!70},
    arr/.style={-Stealth, semithick, draw=black!50},
    splitarr/.style={-Stealth, semithick, draw=black!50}
  ]

  % === SHARED TOP ===
  \node[shared] (agents) {Agent Trajectories};
  \node[desc, below=0.04cm of agents] (agents-d) {Vicsek alignment $+$ Morse attraction--repulsion};

  \node[shared, below=0.35cm of agents-d] (kde) {KDE Density Fields};
  \node[desc, below=0.04cm of kde] (kde-d) {$\rho(x,t)$ on periodic grid via Gaussian kernel};

  \draw[arr] (agents) -- (kde);

  % === BRANCH POINT ===
  \coordinate[below=0.55cm of kde-d] (split);

  % === EF-ROM BRANCH (LEFT) ===
  \node[branchlabel, anchor=south] at ([xshift=-2.9cm]split)
    (efrom-label) {\color{efromLabel}EF-ROM};
  \node[desc, below=0.0cm of efrom-label] (efrom-sub) {Part~I: Data-driven latent model};

  \node[efrom, below=0.25cm of efrom-sub] (align) {\textbf{Shift-Alignment}};
  \node[desc, below=0.04cm of align] (align-d) {FFT cross-correlation on periodic torus};

  \node[efrom, below=0.25cm of align-d] (pod) {\textbf{POD Compression}};
  \node[desc, below=0.04cm of pod] (pod-d) {$\sqrt{\cdot}\,$-transform, truncate to $r$ modes};

  % Place MVAR and LSTM symmetrically about the EF-ROM column centre
  \coordinate[below=0.55cm of pod-d] (stepcentre);
  \node[efrom, minimum width=2.4cm] at ([xshift=-1.4cm]stepcentre) (mvar) {\textbf{MVAR}};
  \node[efrom, minimum width=2.4cm] at ([xshift=1.4cm]stepcentre)  (lstm) {\textbf{LSTM}};
  \node[desc, below=0.04cm of stepcentre, anchor=north] (step-d)
    {Latent time-stepping};

  \node[efrom, below=0.25cm of step-d] (lift) {\textbf{Lift to Density}};
  \node[desc, below=0.04cm of lift] (lift-d) {Inverse POD $+$ undo shift};

  % EF-ROM arrows
  \draw[splitarr] (split) -| (align.north);
  \draw[arr] (align) -- (pod);
  \draw[arr] (pod.south) -- ++(0,-0.65cm) -| (mvar.north);
  \draw[arr] (pod.south) -- ++(0,-0.65cm) -| (lstm.north);
  \draw[arr] (mvar.south) |- ([yshift=0.15cm]lift.west) -- (lift.west);
  \draw[arr] (lstm.south) |- ([yshift=0.15cm]lift.east) -- (lift.east);

  % === WSINDY BRANCH (RIGHT) ===
  \node[branchlabel, anchor=south] at ([xshift=2.9cm]split)
    (ws-label) {\color{wsLabel}WSINDy};
  \node[desc, below=0.0cm of ws-label] (ws-sub) {Part~II: PDE discovery};

  \node[wsindy, below=0.25cm of ws-sub] (fields) {\textbf{Multi-field Extraction}};
  \node[desc, below=0.04cm of fields] (fields-d) {$\rho,\; p_x,\; p_y,\; \Phi = W\ast\rho$};

  \node[wsindy, below=0.25cm of fields-d] (lib) {\textbf{Weak-form Library}};
  \node[desc, below=0.04cm of lib] (lib-d) {Test functions $+$ candidate PDE terms};

  \node[wsindy, below=0.25cm of lib-d] (sparse) {\textbf{Sparse Regression}};
  \node[desc, below=0.04cm of sparse] (sparse-d) {MSTLS with dominant-balance criterion};

  \node[wsindy, below=0.25cm of sparse-d] (etdrk) {\textbf{ETDRK4 Integration}};
  \node[desc, below=0.04cm of etdrk] (etdrk-d) {Fourier-space time-stepping};

  % WSINDy arrows
  \draw[splitarr] (split) -| (fields.north);
  \draw[arr] (fields) -- (lib);
  \draw[arr] (lib) -- (sparse);
  \draw[arr] (sparse) -- (etdrk);

  % === SHARED BOTTOM ===
  % Place evaluation node centered below the two branches
  \coordinate (merge) at ($(lift-d.south)!0.5!(etdrk-d.south) + (0,-0.6cm)$);
  \node[shared] at (merge) (eval) {Evaluation};
  \node[desc, below=0.04cm of eval] (eval-d)
    {RMSE, $R^2$, forecast horizon~$\tau$, mass drift, runtime};

  \draw[splitarr] (lift-d.south) |- ([yshift=0.35cm]eval.west) -- (eval.west);
  \draw[splitarr] (etdrk-d.south) |- ([yshift=0.35cm]eval.east) -- (eval.east);

\end{tikzpicture}
\caption{Overview of the two complementary modeling pipelines developed in this
  thesis.  Both branches share the same microscopic-to-macroscopic data
  generation (top) and are evaluated on common metrics (bottom).
  Left: the equation-free reduced-order model compresses aligned density
  snapshots into a low-dimensional latent space and evolves them with learned
  time-steppers.
  Right: WSINDy discovers interpretable PDE flux laws from multi-field data via
  weak-form sparse regression.}
\label{fig:pipeline_overview}
\end{figure}
